% ============================================================
% CHAPITRE 6 (Sprint 4)
% ============================================================
\chapter{Sprint 4 : Back-Office et Analytics}

\section{Introduction}

Le Sprint 4 marque la transition vers l'administration globale de la plateforme Assiette Gala. Si les sprints précédents se sont focalisés sur la construction de l'expérience d'achat (vitrine, commandes, événements), ce sprint vise à fournir à la direction les outils de pilotage stratégique de l'entreprise. 

Ce cycle consolide l'ensemble du "Back-Office" : le tableau de bord analytique, le contrôle des flux financiers (dépenses, encaissements), le pilotage des stocks et des fournisseurs, ainsi que la génération des documents comptables légaux (factures, bons de commande). Le but est de remplacer tous les registres papier par un système d'information centralisé.

% --------------------------------------------------
\section{Backlog du sprint}

Le Tableau~\ref{tab:backlog_sprint4_admin} recense les user stories priorisées pour ce sprint.

\begin{table}[H]
    \centering
    \caption{Backlog du sprint 4 : Back-Office et Analytics}
    \label{tab:backlog_sprint4_admin}
    \renewcommand{\arraystretch}{1.3}
    \begin{tabular}{|c|p{8.5cm}|c|}
        \hline
        \textbf{ID} & \textbf{User Story} & \textbf{Durée (jours)} \\
        \hline
        US-18 & En tant qu'administrateur, je souhaite visualiser des KPI et graphiques d'évolution afin de piloter l'activité. & 4 \\
        \hline
        US-19 & En tant qu'administrateur, je souhaite analyser la répartition des commandes par statut et les tendances. & 2 \\
        \hline
        US-20 & En tant qu'administrateur, je souhaite suivre les flux financiers (caisse, dépenses) pour garantir la rentabilité. & 3 \\
        \hline
        US-21 & En tant qu'administrateur, je souhaite générer des documents comptables au format PDF. & 3 \\
        \hline
    \end{tabular}
\end{table}

% --------------------------------------------------
\section{Partie A : Tableau de Bord et Analytics}

\subsection{Architecture du système de reporting}

Le module d'analytics repose sur une architecture en trois couches :
\begin{enumerate}
    \item \textbf{Base de données (Prisma) :} Utilisation de la fonction \texttt{groupBy} et des agrégations SQL pour consolider les données côté serveur sans surcharger la mémoire Node.js.
    \item \textbf{API Gateway :} Un contrôleur dédié aux statistiques orchestre les appels concurrents via \texttt{Promise.all()} pour minimiser le temps de réponse.
    \item \textbf{Frontend (Recharts) :} Les données brutes sont transformées et injectées dans des composants graphiques interactifs (SVG) fluides.
\end{enumerate}

\subsection{Indicateurs de Performance (KPI)}

Le tableau de bord supérieur offre une vue synthétique des métriques vitales :
\begin{itemize}
    \item Chiffre d'affaires total (avec comparaison par rapport au mois précédent).
    \item Nombre de commandes et taux de validation.
    \item Panier moyen.
    \item Volume d'événements en attente de traitement.
\end{itemize}

\subsection{Visualisations Graphiques}

L'interface analytique intègre plusieurs graphiques interactifs :
\begin{itemize}
    \item \textbf{Courbe d'évolution des revenus :} Permet d'identifier la saisonnalité de l'activité.
    \item \textbf{Diagramme en beignet (Donut Chart) :} Répartition des commandes par statut (En attente, En préparation, Livré).
    \item \textbf{Graphique en barres :} Palmarès des plats les plus vendus (Top 5).
\end{itemize}

% --------------------------------------------------
\section{Partie B : Suivi des Flux Financiers}

Le module finance assure le suivi complet des mouvements monétaires, consolidant les ventes issues du module de commandes et les charges de fonctionnement :

\begin{itemize}
    \item \textbf{Encaissements} : Visualisation des entrées d'argent par mode de paiement, automatiquement alimentée par le changement de statut des commandes (passage en "Complété").
    \item \textbf{Dépenses (Charges) :} Interface permettant l'enregistrement des sorties financières (factures fournisseurs, achats de matières premières, charges fixes) avec catégorisation.
    \item \textbf{Balance} : Calcul du flux net de trésorerie sur la période sélectionnée.
\end{itemize}

% --------------------------------------------------
\section{Partie C : Gestion Documentaire (PDF)}

La génération de documents légaux constitue la dernière brique de l'ERP métier. Le module documentaire permet la création à la volée de :

\begin{itemize}
    \item \textbf{Factures (Invoices) :} Génération dynamique via le package \texttt{pdfkit}. Les factures intègrent le logo d'Assiette Gala, les mentions légales tunisiennes, la numérotation séquentielle inviolable (format \texttt{INV-YYYYMM-<hex>}), le détail des lignes de commande et le calcul de la TVA à 19\%.
    \item \textbf{Bons de commande fournisseurs :} Documents d'achat liant les besoins en stock aux fournisseurs enregistrés.
\end{itemize}

% --------------------------------------------------
\section{Conception}

\subsection{Diagramme de cas d'utilisation}

\begin{figure}[H]
    \centering
    \includegraphics[width=0.85\textwidth]{uc_sprint4.png}
    \caption{Diagramme de cas d'utilisation : Back-Office et Analytics}
    \label{fig:uc_sprint4}
\end{figure}

\subsection{Description textuelle des cas d'utilisation}

\subsubsection{Cas d'utilisation : Visualiser les KPI}

\begin{table}[H]
    \centering
    \caption{Description textuelle : Visualiser les KPI}
    \label{tab:desc_kpi}
    \renewcommand{\arraystretch}{1.3}
    \begin{tabular}{|p{4cm}|p{8.5cm}|}
        \hline
        \textbf{Cas d'utilisation} & Visualiser les KPI et graphiques \\
        \hline
        \textbf{Acteur principal} & Administrateur \\
        \hline
        \textbf{Précondition} & Authentifié avec le rôle \texttt{ADMIN} \\
        \hline
        \textbf{Postcondition} & Les indicateurs sont affichés à jour \\
        \hline
        \textbf{Scénario principal} &
        \begin{enumerate}
            \item L'administrateur accède à la page d'accueil du Dashboard
            \item Le système interroge l'API Gateway (\texttt{/api/stats/overview})
            \item L'API consolide les données via Prisma (\texttt{groupBy})
            \item Affichage des cartes animées (CA, Commandes)
            \item Rendu des graphiques Recharts (Courbe de revenus, Donut des statuts)
        \end{enumerate} \\
        \hline
    \end{tabular}
\end{table}

\subsubsection{Cas d'utilisation : Suivre les flux financiers}

\begin{table}[H]
    \centering
    \caption{Description textuelle : Flux financiers}
    \label{tab:desc_finance}
    \renewcommand{\arraystretch}{1.3}
    \begin{tabular}{|p{4cm}|p{8.5cm}|}
        \hline
        \textbf{Cas d'utilisation} & Enregistrer une dépense et suivre la caisse \\
        \hline
        \textbf{Acteur principal} & Administrateur \\
        \hline
        \textbf{Précondition} & Authentifié avec le rôle \texttt{ADMIN} \\
        \hline
        \textbf{Postcondition} & La dépense est enregistrée, la balance financière est mise à jour \\
        \hline
        \textbf{Scénario principal} &
        \begin{enumerate}
            \item Navigation vers le module "Finance"
            \item Saisie d'une facture fournisseur (ex: 500 TND pour viandes)
            \item Validation du formulaire
            \item Le système déduit le montant des encaissements globaux
            \item Inscription de la transaction dans le journal d'audit
        \end{enumerate} \\
        \hline
    \end{tabular}
\end{table}

\subsubsection{Cas d'utilisation : Générer un document PDF}

\begin{table}[H]
    \centering
    \caption{Description textuelle : Documents comptables}
    \label{tab:desc_pdf}
    \renewcommand{\arraystretch}{1.3}
    \begin{tabular}{|p{4cm}|p{8.5cm}|}
        \hline
        \textbf{Cas d'utilisation} & Générer une facture au format PDF \\
        \hline
        \textbf{Acteur principal} & Administrateur / Client \\
        \hline
        \textbf{Précondition} & Commande clôturée (\texttt{COMPLETED}) \\
        \hline
        \textbf{Postcondition} & Fichier PDF téléchargé \\
        \hline
        \textbf{Scénario principal} &
        \begin{enumerate}
            \item Clic sur le bouton d'export PDF de la commande
            \item Le module \texttt{pdfkit} dessine l'en-tête (Logo Assiette Gala, Matricule)
            \item Itération sur les lignes de la commande pour tracer le tableau
            \item Calcul de la TVA et du TTC
            \item Envoi du flux binaire en téléchargement
        \end{enumerate} \\
        \hline
    \end{tabular}
\end{table}

\subsection{Diagrammes de séquence}

\subsubsection{Chargement du tableau de bord}

\begin{figure}[H]
    \centering
    \includegraphics[width=0.9\textwidth]{seq_dashboard_load.png}
    \caption{Diagramme de séquence : Chargement du Dashboard}
    \label{fig:seq_dashboard}
\end{figure}

\subsubsection{Génération d'une facture PDF}

\begin{figure}[H]
    \centering
    \includegraphics[width=0.9\textwidth]{seq_facture.png}
    \caption{Diagramme de séquence : Génération de facture PDF}
    \label{fig:seq_facture}
\end{figure}

\subsubsection{Filtrage des statistiques par période}

\begin{figure}[H]
    \centering
    \includegraphics[width=0.9\textwidth]{seq_stats_filter.png}
    \caption{Diagramme de séquence : Filtrage des statistiques par période}
    \label{fig:seq_stats_filter}
\end{figure}

\subsection{Diagramme de classes}

\begin{figure}[H]
    \centering
    \includegraphics[width=0.95\textwidth]{class_sprint4.png}
    \caption{Diagramme de classes du sprint 4 : Back-Office et Analytics}
    \label{fig:class_sprint4}
\end{figure}
% --------------------------------------------------
\section{Réalisation technique}

\subsection{Points d'entrée de l'API}

\begin{table}[H]
    \centering
    \caption{Endpoints REST du Back-Office}
    \label{tab:endpoints_admin}
    \renewcommand{\arraystretch}{1.2}
    \begin{tabular}{|l|p{4.5cm}|p{5.5cm}|}
        \hline
        \textbf{Méthode} & \textbf{Route} & \textbf{Description} \\
        \hline
        GET & /api/stats/overview & Récupération des KPI globaux \\
        \hline
        GET & /api/stats/revenue-chart & Données pour la courbe des revenus \\
        \hline
        GET & /api/finance/transactions & Liste des flux (entrées/sorties) \\
        \hline
        POST & /api/finance/expense & Enregistrement d'une dépense \\
        \hline
        GET & /api/documents/invoice/:id & Génération de facture PDF \\
        \hline
    \end{tabular}
\end{table}

\subsection{Tests fonctionnels}

\begin{table}[H]
    \centering
    \caption{Tests fonctionnels : Sprint 4}
    \label{tab:tests_sprint4}
    \renewcommand{\arraystretch}{1.3}
    \begin{tabular}{|c|p{4cm}|p{4cm}|p{2cm}|}
        \hline
        \textbf{N°} & \textbf{Cas de test} & \textbf{Comportement attendu} & \textbf{Résultat} \\
        \hline
        1 & Accès aux KPI sans token admin & Refus (401 ou 403) & Réussi \\
        \hline
        2 & Calcul du chiffre d'affaires & Somme exacte des commandes livrées & Réussi \\
        \hline
        3 & Enregistrement d'une dépense & Dépense créée et déduite de la balance & Réussi \\
        \hline
        4 & Génération facture PDF & PDF généré avec TVA à 19\% et lignes correctes & Réussi \\
        \hline
    \end{tabular}
\end{table}

% --------------------------------------------------
\section*{Conclusion}
\addcontentsline{toc}{section}{Conclusion}

\begin{samepage}
Ce Sprint 4 transforme la plateforme d'une simple interface de prise de commande en un véritable outil de pilotage (\textbf{ERP allégé}). La direction d'Assiette Gala dispose désormais d'une visibilité totale sur ses indicateurs de croissance, sa trésorerie et l'édition de ses documents comptables.

L'administration étant désormais entièrement digitalisée et outillée, le système est prêt à accueillir le cinquième et dernier sprint : l'intégration de la couche algorithmique d'Intelligence Artificielle.
\end{samepage}

